\chapter{Exemples de fichiers manifest}
\label{annexe:manifests}

Les fichiers manifest constituent l'artefact intermédiaire central du framework ETL.
Ils encodent la totalité des règles de mapping, de résolution de clés étrangères et de
logique métier d'un import sous une forme lisible par Claude Code et transformable en
pipelines Apache Hop exécutables. Deux exemples sont présentés ci-dessous pour illustrer
la complexité croissante du format.

Le manifest \texttt{partner.json} présenté ci-dessous couvre l'import de tiers avec
adresses embarquées, tables de jonction et post-traitement SQL — l'un des cas les plus
complets du framework.

% ──────────────────────────────────────────────────────────────
\section{\texttt{partner.json} — Import des tiers}
\label{annexe:manifest-partner}
% ──────────────────────────────────────────────────────────────

\lstinputlisting[
  language=json,
  caption={Manifest d'import des tiers (\texttt{partner.json})},
  label={lst:manifest-partner}
]{../../hop-pfe-demo/v2/manifests/partner.json}

% ══════════════════════════════════════════════════════════════
\chapter{Structure du skill \texttt{analyze-and-map-entity}}
\label{annexe:skill}
% ══════════════════════════════════════════════════════════════

Le skill \texttt{/hop:analyze-and-map-entity} est l'agent de mapping conversationnel du
framework. Invocable depuis Claude Code ou depuis l'interface Axelor Import Studio, il
lit un fichier CSV, mappe ses colonnes vers une table Axelor cible, résout chaque
dépendance de clé étrangère en interrogeant la base de données et en dialoguant avec
le consultant, puis écrit le fichier \texttt{manifest.json} résultant.

\medskip
\textbf{Fichiers de référence lus avant la classification} :
\begin{itemize}
  \item \texttt{common-entities.md} — connaissance des entités product, partner, sale\_order, invoice et leurs cibles FK
  \item \texttt{fk-reference.md} — correspondances nom de colonne $\to$ ref\_table, stratégies de lookup, valeurs fk\_source
  \item \texttt{column-types.md} — règles de typage, colonnes exclues, conventions manifest (dtype, change\_detection, audit\_defaults)
  \item \texttt{manifest.example.json} — bibliothèque de quatre manifests v2.3 complets et validés
  \item \texttt{address-parsing.md} — lu uniquement si le CSV contient une adresse postale en texte libre
  \item \texttt{advanced-entities.md} — lu uniquement pour les entités multi-fichiers (en-tête + lignes + jonction)
\end{itemize}

\bigskip

% ── Phase 0 ──────────────────────────────────────────────────
\section*{Phase 0 — Vérification de la connexion DB}

Avant toute action, teste la connexion PostgreSQL via \texttt{psql}. Si la connexion
réussit, les vérifications FK seront en direct ; en cas d'échec, le skill continue en
mode dégradé (résolution par réponses utilisateur uniquement).

% ── Phase 1 ──────────────────────────────────────────────────
\section*{Phase 1 — Lecture du CSV et du schéma DB}

Les deux opérations sont lancées simultanément :

\begin{itemize}
  \item Lecture du CSV (en-têtes + $\sim$20 premières lignes pour échantillonnage des valeurs)
  \item Requêtes \texttt{information\_schema} sur la table cible — colonnes réelles + contraintes FK
\end{itemize}

\noindent\textbf{Phase 1.2 — Clé d'import (\texttt{import\_id})} : identifie la colonne
servant de clé de dédup upsert. Si \texttt{importId} est présent, l'utilise directement.
Sinon, propose de construire la clé depuis \texttt{id} avec un préfixe configurable
(ex. : \texttt{DOLI\_PROD\_12345}).

\medskip
\noindent\textbf{Phase 1.3 — Classification unifiée des colonnes} : chaque colonne est
scorée sur trois signaux simultanés :

\begin{center}
\renewcommand{\arraystretch}{1.4}
{\small
\begin{tabular}{|p{3.4cm}|p{4.0cm}|p{6.0cm}|}
  \hline
  \textbf{Signal} & \textbf{Exemple} & \textbf{Interprétation} \\
  \hline
  A — Notation pointée & \texttt{productFamily.}\newline\texttt{importId} & FK certain — table et stratégie explicites \\
  \hline
  B — Motif du nom & \texttt{currencyCode} & FK probable — suffixe \texttt{Code} / \texttt{Id} / \texttt{Name} \\
  \hline
  C — Échantillon de valeurs & faible cardinalité, valeurs courtes & FK corroborant \\
  \hline
\end{tabular}}
\end{center}

\noindent La combinaison des signaux produit quatre classifications :
\textbf{FK certain} (traité sans question), \textbf{FK probable} (confirmation rapide),
\textbf{Incertain} (présentation des preuves + question), \textbf{Scalaire} (ajouté
directement à \texttt{column\_mappings}).

% ── Phase 2 ──────────────────────────────────────────────────
\section*{Phase 2 — Vérification DB de chaque FK}

Pour chaque FK identifiée, interroge la base pour savoir si les valeurs CSV existent déjà :

\begin{lstlisting}[language={}, basicstyle=\ttfamily\footnotesize, frame=single,
                   backgroundcolor=\color{gray!4}]
SELECT DISTINCT <lookup_col> FROM <ref_table>
WHERE <lookup_col> IN (<valeurs_csv>) LIMIT 50;
\end{lstlisting}

\noindent Les FKs pointant vers la même \texttt{ref\_table} sont regroupées en une seule
requête. Trois résultats possibles :

\begin{itemize}
  \item \textbf{[OK] Toutes trouvées} $\to$ \texttt{fk\_source: preloaded}, aucune question
  \item \textbf{[WARNING] Partiellement trouvées} $\to$ question sur les valeurs manquantes (Phase 3)
  \item \textbf{[MISSING] Aucune trouvée} $\to$ question sur la stratégie complète (Phase 3)
\end{itemize}

% ── Phase 3 ──────────────────────────────────────────────────
\section*{Phase 3 — Résolution des FK manquantes}

Pour chaque \texttt{ref\_table} avec des valeurs absentes, le skill pose une question
structurée via \texttt{AskUserQuestion} et génère la configuration correspondante :

\begin{center}
\renewcommand{\arraystretch}{1.3}
\begin{tabular}{|p{3.5cm}|p{3.5cm}|p{6.2cm}|}
  \hline
  \textbf{Choix utilisateur} & \textbf{fk\_source} & \textbf{Ce qui est généré} \\
  \hline
  Fichier séparé & \texttt{separate\_entity} & Section \texttt{entities\{\}} complète + dépendance FK \\
  \hline
  Créer automatiquement & \texttt{embedded} & Pre-pipeline \texttt{create\_missing\_<ref>} + \texttt{pre\_pipeline\_config} \\
  \hline
  Valeur par défaut & \texttt{constant} & \texttt{constant\_value} vérifié en DB + \texttt{on\_miss: block\_row} \\
  \hline
  Ignorer & \texttt{set\_null} & \texttt{on\_miss: set\_null}, aucun pipeline additionnel \\
  \hline
\end{tabular}
\end{center}

\noindent\textbf{Cas spéciaux détectés automatiquement} :
\begin{itemize}
  \item \textbf{Adresse postale embarquée} — colonne de type \texttt{"70 RUE PASTEUR | 06000 NICE | FRANCE"} :
        déclenche la lecture de \texttt{address-parsing.md} et génère une entité
        \texttt{address} séparée avec \texttt{parse\_config} (RegexEval + dédup).
  \item \textbf{Entité multi-fichiers} — CSV en-tête + CSV lignes + CSV jonction :
        déclenche \texttt{advanced-entities.md} et génère
        \texttt{junction\_pipelines[]}, \texttt{post\_pipelines[]},
        \texttt{status\_sequence} et \texttt{workflow\_sql\_action}.
\end{itemize}

% ── Phase 4 ──────────────────────────────────────────────────
\section*{Phase 4 — Point de confirmation}

\textbf{Phase 4.1 — Vérification des FK obligatoires} : croise la liste des FK
obligatoires (fournie par \texttt{registry.py}) avec les FK résolues. Toute FK
obligatoire sans colonne CSV correspondante déclenche une question sur l'action à prendre
(fichier séparé / valeur par défaut / laisser NULL avec avertissement).

\medskip
\textbf{Phase 4.2 — Tableau récapitulatif} : présente l'ensemble des résolutions avant
écriture :

\begin{center}
\renewcommand{\arraystretch}{1.35}
{\small
\begin{tabular}{|p{3.3cm}|p{3.2cm}|p{2.7cm}|p{4.0cm}|}
  \hline
  \textbf{ref\_table} & \textbf{Colonnes CSV} & \textbf{Stratégie} & \textbf{Détails} \\
  \hline
  \texttt{base\_currency}        & \texttt{saleCurrency.}\newline\texttt{code}     & preloaded       & EUR, USD trouvés [OK] \\
  \hline
  \texttt{base\_unit}            & \texttt{unit}                  & embedded        & 2/4 manquants $\to$ pre-pipeline \\
  \hline
  \texttt{base\_product\_}\newline\texttt{family} & \texttt{productFamily.}\newline\texttt{importId}& separate\_entity& depends\_on: product\_family \\
  \hline
  \texttt{base\_company}         & --                             & constant        & AXE (toutes les lignes) \\
  \hline
\end{tabular}}
\end{center}

% ── Phase 5 ──────────────────────────────────────────────────
\section*{Phase 5 — Écriture du manifest}

Récupère le timestamp réel via \texttt{date -Iseconds}, résout le chemin absolu de
sortie, puis écrit \texttt{manifest.json} avec l'outil \texttt{Write} en suivant
le schéma v2.3 complet :

\begin{lstlisting}[language={}, basicstyle=\ttfamily\footnotesize, frame=single,
                   backgroundcolor=\color{gray!4}]
{
  "version": "2.3",
  "generated_at": "<ISO timestamp>",
  "generated_by": "axelor-mapper v2 (claude agent)",
  "primary_entity": "<entity_key>",
  "import_order": ["<dep1>", "<dep2>", "<primary>"],
  "entities": { ... },          // une section par entite
  "junction_pipelines": [ ... ],// M2M uniquement
  "post_pipelines": [ ... ]     // calculs post-import
}
\end{lstlisting}

\noindent Une fois le manifest écrit, le consultant peut lancer le skill
\texttt{/hop:generate-pipeline} avec ce manifest pour produire les pipelines
Apache Hop correspondants.

